\section{The Architect's Question}\label{the-architects-question}

After this chapter we should be able to reason about the two fundamental
modes of LLM inference --- prefill and decode --- and to identify which
system constraint each mode exposes. We should be able to compute the
bandwidth and FLOP demands of a given workload, compare them against
real hardware limits, and understand why decode is HBM-bandwidth-bound
while prefill is compute-bound. This distinction directly shapes
architectural decisions: whether we need more GPUs, whether we can
disaggregate prefill from decode, and how we size memory and compute in
a fleet. After this chapter we should see inference not as a single
monolithic cost but as two opposite-bound processes that require
different levers.

\section{Concept}\label{concept}

Inference is the process of turning a sequence of input tokens into a
sequence of output tokens using a trained neural network. In an
autoregressive model, every generated token requires a fresh forward
pass, but the computation split between the prefix (the prompt) and the
suffix (the generated output) is fundamentally different.

\textbf{Prefill} is the first pass over the prompt. The model reads
every token of the input context, computes attention over the entire
prefix, and produces the first output token. During prefill, the model
attends over all input positions, so the work grows with the square of
the context length.

\textbf{Decode} is the per-token generation that follows. After prefill,
each new token only attends to the cached keys and values from the
prefix plus all previously generated tokens. In the
\emph{batch-1-equivalent} weight-streaming model we use for this
first-pass arithmetic, the per-step cost is constant in context length
--- we only attend to the new token against the cache --- and it
requires reading the entire model weight matrix from HBM every step,
because the auto-regressive pass reuses the same parameters over and
over. (In a real serving stack, batching and weight-read amortization
change this per-request picture substantially; that is the subject of
later chapters.)

The critical architectural insight: prefill is dominated by FLOPs
(matrix-multiply arithmetic), while decode is dominated by HBM memory
bandwidth (reading 140 GB of weights 40 times a second). These are
opposite bottlenecks, and confusing them is the most common
architectural misstep.

\begin{figure}
\centering
\pandocbounded{\includegraphics[keepaspectratio,alt={The canonical end-to-end inference pipeline: prompt to generated tokens, divided into PREFILL (top) and DECODE (bottom). Prefill is the one-shot parallel burst that populates the KV cache and yields the first token --- it stresses compute and is measured by TTFT. Decode is the sequential autoregressive loop that reuses the cached K/V and re-reads weights every step --- it stresses HBM bandwidth and capacity, and is measured by TPOT/ITL. {[}ILLUSTRATIVE conceptual{]}},width=\maxwidth{\textwidth},keepaspectratio]{figures/fig-02-0202.pdf}}
\caption{The canonical end-to-end inference pipeline: prompt to generated tokens, divided into PREFILL (top) and DECODE (bottom). Prefill is the one-shot parallel burst that populates the KV cache and yields the first token --- it stresses compute and is measured by TTFT. Decode is the sequential autoregressive loop that reuses the cached K/V and re-reads weights every step --- it stresses HBM bandwidth and capacity, and is measured by TPOT/ITL. {[}ILLUSTRATIVE conceptual{]}}

\label{fig:2.1}\end{figure}

\emph{This is the figure the rest of the book builds on.} Keep the shape
in mind: every later chapter refines one part of it. \textbf{\hyperref[chap:3]{Chapter 3}}
(model) asks \emph{what the blocks compute and how the KV cache is
sized}; \textbf{6} (metrics) measures which phase is the bottleneck;
\textbf{7} (memory) prices the KV cache precisely; \textbf{8} (compute)
derives the FLOPs and roofline that make prefill compute-bound;
\textbf{10} (parallelism) splits the model across the devices the
pipeline traverses; \textbf{11} (serving) schedules many of these
pipelines; \textbf{15} (performance) diagnoses which phase actually
hurts; and the fleet chapters price the host count the loop demands.
Because the two phases stress \emph{different} resources, the armature
of the whole book sits on this single split.

\textbf{Where the Q/K/V and KV cache live in this picture (the bridge to
Ch. 3/7).} The transformer layer computes, for each token, three learned
projections --- \textbf{query (Q)}, \textbf{key (K)} and \textbf{value
(V)}. Attention scores are the similarity of a query to every prior key;
the scores weight the values, which are blended into the token's new
representation. What matters for architecture: \textbf{K and V for a
token do not change once that token has been processed.} The query for a
\emph{later} token attends against them, but rewriting them would be
redundant work. So decode does not recompute the prefix's K/V --- it
\emph{reuses} the stored K and V from prefill. That is the KV cache: not
an optimization bolted on, but the direct consequence of \emph{which}
projections are reusable and which (the query) is always new. It is also
why KV memory grows with sequence length (one K and one V per token, per
layer) and why attention implementation and precision (GQA/MQA,
FlashAttention, KV quantization) change an architecture problem that is
really a memory problem. Ch. 3 and Ch. 7 do the arithmetic; this figure
is where the causal thread starts.

\textbf{A vocabulary note that prevents a common conflation.} Several
terms sound like the same thing but are \emph{not} interchangeable, and
the chapter will use them precisely:

\begin{longtable}[]{@{}
  >{\raggedright\arraybackslash}p{0.2933\linewidth}
  >{\raggedright\arraybackslash}p{0.2933\linewidth}
  >{\raggedright\arraybackslash}p{0.2933\linewidth}@{}}
\toprule\noalign{}
\begin{minipage}[b]{\linewidth}\raggedright
term
\end{minipage} & \begin{minipage}[b]{\linewidth}\raggedright
what it is
\end{minipage} & \begin{minipage}[b]{\linewidth}\raggedright
what it is \textbf{not}
\end{minipage} \\
\midrule\noalign{}
\endhead
\bottomrule\noalign{}
\endlastfoot
\textbf{KV cache} & the stored K and V tensors per token per layer,
saved so decode does not recompute the prefix & an \emph{optimization}
over a shared prefix \\
\textbf{PagedAttention} & a memory-\emph{allocation} scheme for the KV
cache (splits it into fixed-size pages, cuts fragmentation/waste) & a
way to \emph{reduce} the amount of KV state \\
\textbf{prefix caching / KV reuse} & reusing \emph{already-computed} K/V
for a shared prefix, to skip re-prefilling it & the same as caching the
KV of a \emph{single} request's own decode \\
\textbf{prompt caching} & broader system/provider term; the book uses it
to mean \emph{reusing the KV of a shared prompt prefix across requests}
& a distinct mechanism from prefix caching \\
\textbf{KV quantization} & reducing the \emph{bytes} per K/V element
(e.g.~FP16→FP8) & reducing the \emph{number} of key/value entries \\
\textbf{KV offloading} & moving some KV state out of GPU HBM (to host
DRAM/SSD), trading capacity against transfer latency & shrinking the
state itself \\

\label{tab:2.1}\end{longtable}

The short version: the \textbf{KV cache} is \emph{the state};
\textbf{PagedAttention} manages \emph{where that state lives} on the
device; \textbf{prefix caching} avoids \emph{computing it twice} for a
shared prefix; \textbf{quantization} shrinks \emph{each element};
\textbf{offloading} moves \emph{some of it elsewhere}. Each attacks a
different constraint (capacity, fragmentation, redundant compute, byte
count, residency), and each has its own trade-off --- later chapters put
numbers on them.

\textbf{The resource model (the book's recurring framework).} Inference
consumes several \emph{fundamentally different} resources, and the
single most useful habit an architect can build is to name which one a
given optimization or bottleneck is about. The pipeline figure above
surfaces four of them, and the metric chapter (Ch. 6) and the design
chapter (Ch. 12) build on the same set. The book will use these names
consistently:

\begin{longtable}[]{@{}
  >{\raggedright\arraybackslash}p{0.2200\linewidth}
  >{\raggedright\arraybackslash}p{0.2200\linewidth}
  >{\raggedright\arraybackslash}p{0.2200\linewidth}
  >{\raggedright\arraybackslash}p{0.2200\linewidth}@{}}
\toprule\noalign{}
\begin{minipage}[b]{\linewidth}\raggedright
resource
\end{minipage} & \begin{minipage}[b]{\linewidth}\raggedright
unit
\end{minipage} & \begin{minipage}[b]{\linewidth}\raggedright
what it limits
\end{minipage} & \begin{minipage}[b]{\linewidth}\raggedright
first surfaces in
\end{minipage} \\
\midrule\noalign{}
\endhead
\bottomrule\noalign{}
\endlastfoot
\textbf{compute capacity} & FLOPs & how much arithmetic per unit time &
prefill (Ch. 2, 8) \\
\textbf{memory bandwidth} & bytes/s & how fast data can move & decode
(Ch. 2, 8) \\
\textbf{memory capacity} & bytes resident & how much state fits (weights
+ KV) & KV cache / concurrency (Ch. 7) \\
\textbf{interconnect bandwidth/latency} & bytes/s, delay & cross-GPU /
cross-host movement & parallelism (Ch. 9, 10) \\
\textbf{scheduling capacity} & useful work resident & keeping hardware
busy despite request raggedness & serving (Ch. 11) \\
\textbf{latency budget} & TTFT, TPOT/ITL, end-to-end & the SLO the
system must meet & workload (Ch. 4) \\
\textbf{throughput / goodput} & useful tokens/s under SLO & completed
work, not just tokens/s & metrics (Ch. 6) \\

\label{tab:2.2}\end{longtable}

The discipline the book applies everywhere: \textbf{every optimization
is a hypothesis about which resource it shifts and what it trades away.}
FlashAttention moves the bottleneck from \emph{bandwidth} toward
\emph{compute} (fewer HBM round-trips); KV quantization trades
\emph{capacity} for a small precision loss; prefix caching trades
\emph{prefill compute} for \emph{cache residency}; continuous batching
trades \emph{latency} for \emph{utilization}; P/D disaggregation trades
a \emph{homogeneous pool's} simplicity for separate \emph{compute- vs
bandwidth-optimized} pools. When the book introduces a technique, it
will name the resource it targets and the resource it costs --- that is
the difference between reasoning and a vendor catalog. The optimization
decision map (Ch. 15) is this framework read backwards: from an observed
symptom to the resource that is likely starved.

\textbf{The causal chain (how the chapters hold together).} The book is
one argument, not a set of topics. The thread is:

tokens determine sequence length → sequence length determines KV growth
and attention work → KV growth determines memory pressure → attention
and model execution determine compute and bandwidth demand → those
resource demands determine latency behavior → latency and concurrency
determine scheduler requirements → scheduler behavior determines serving
efficiency → serving efficiency and SLOs determine host count → host
count and model topology determine interconnect and fleet architecture →
fleet architecture determines cost, resilience, placement, and
operational complexity.

Every chapter is one link of that chain, and each \emph{deepens} the
same mental model from Ch. 2 rather than restarting it. When you reach
Ch. 11 (serving) or Ch. 17 (fleet), the concepts were introduced in Ch.
2 --- Ch. 2 gave the shape, the later chapters give the numbers and the
decisions. If a chapter ever feels like ``here is another concept,'' it
has broken the thread.

\section{Mental Model}\label{mental-model}

Think of inference as two fundamentally different physical processes,
distinguished by what limits them.

\textbf{Prefill is compute-bound.} It is like turning on a massive
industrial fire hose to fill a reservoir. The burst of flow is enormous,
and the time it takes is limited entirely by the raw horsepower of the
pump (compute/FLOPs) pushing the water. The pipe is wide open, but the
motor is working at its absolute limit.

\textbf{Decode is bandwidth-bound.} It is like trying to empty a
140,000-gallon swimming pool through a standard garden hose every 25
milliseconds just to produce a single glass of water. The pump (compute)
is more than powerful enough, but the diameter of the hose (HBM
bandwidth) physically cannot move that volume of water in that time
frame.

When we ask ``is this workload compute-bound or bandwidth-bound?'' the
answer depends entirely on which physical constraint we are hitting: the
horsepower of the engine, or the diameter of the pipe.

\section{Worked Example}\label{worked-example}

To make the distinction concrete, let us walk through the canonical
enterprise Q\&A scenario from canonical scenario (Ch 4, \hyperref[tab:4.3]{Table 4-3}): a
70B-class dense model in FP16, 1 host with 8×H100 GPUs, prompt of 1,200
tokens + 8K retrieved context (\textasciitilde9.2K input), 300-token
output, TTFT budget 1.2 s (retrieval \textasciitilde120 ms + prefill),
TPOT budget \textasciitilde25 ms/token. All numbers are derived from the
canonical scenario; none are measurement claims. Where the book writes
\emph{Analytical bound --- not a benchmark}, it marks a first-order
model check under a stated bounding assumption, not a measurement of
actual serving throughput.

\subsection{Table 2-1 --- Decode bandwidth and prefill FLOPs (worked
example, not
reference)}\label{table-2-1-decode-bandwidth-and-prefill-flops-worked-example-not-reference}

\begin{longtable}[]{@{}
  >{\raggedright\arraybackslash}p{0.2933\linewidth}
  >{\raggedright\arraybackslash}p{0.2933\linewidth}
  >{\raggedright\arraybackslash}p{0.2933\linewidth}@{}}
\toprule\noalign{}
\begin{minipage}[b]{\linewidth}\raggedright
metric
\end{minipage} & \begin{minipage}[b]{\linewidth}\raggedright
value
\end{minipage} & \begin{minipage}[b]{\linewidth}\raggedright
derivation
\end{minipage} \\
\midrule\noalign{}
\endhead
\bottomrule\noalign{}
\endlastfoot
Model params & 70B & canonical scenario (Ch 4, \hyperref[tab:4.3]{Table 4-3}) \\
Model weight bytes (FP16) & 140 GB & 2 bytes × 70B params
{[}ILLUSTRATIVE{]}{[}DERIVED{]} \\
Decode: weight-read per token & 140 GB & Auto-regressive,
batch-1-equivalent: one read of all weights per generated token before
batching/amortization {[}ILLUSTRATIVE{]}{[}DERIVED{]} \\
Decode: required bandwidth (TPOT \textasciitilde25 ms) & 5.6 TB/s & 140
GB / 0.025 s (batch-1-equivalent weight-streaming model, before
amortization) {[}DERIVED{]} \\
H100 HBM3 peak bandwidth & 3.35 TB/s & NVIDIA H100 specs
{[}1P{]}{[}FACT{]} \\
Decode: bandwidth verdict & bandwidth-bound & 5.6 \textgreater{} 3.35 →
under the batch-1-equivalent model, a single H100 cannot meet the
demand; batching/amortization change this (see later chapters)
{[}ILLUSTRATIVE{]}{[}DERIVED{]} \\
Prefill: input tokens & 9,200 & 1,200 prompt + 8K context
{[}ILLUSTRATIVE{]}{[}DERIVED{]} \\
Prefill: FLOPs (2 × params × tokens) & 1.29 PFLOP & 2 ×70e9 ×9.2e3 ≈
1.29 ×10\^{}15 {[}DERIVED{]} \\
H100 BF16 dense compute & 989 TFLOPS & NVIDIA H100 BF16 tensor-core peak
{[}1P{]}{[}FACT{]} \\
Prefill: compute verdict & compute-bound & 1.29 PFLOP /
\textasciitilde1.08 s ≈ 1.19 PFLOPS \textgreater{} 0.989 PFLOPS peak →
single H100 insufficient for real-time prefill {[}DERIVED{]} \\

\label{tab:2.3}\end{longtable}

\begin{figure}
\centering
\pandocbounded{\includegraphics[keepaspectratio,alt={Decode vs prefill: bandwidth vs compute {[}ILLUSTRATIVE conceptual{]}},width=\maxwidth{\textwidth},keepaspectratio]{figures/fig-02-0201.pdf}}
\caption{Decode vs prefill: bandwidth vs compute {[}ILLUSTRATIVE conceptual{]}}

\label{fig:2.2}\end{figure}

\emph{Decode is HBM-bandwidth-bound (5.6 TB/s vs H100 3.35 TB/s);
prefill is compute-bound (\textasciitilde1.19 PFLOPS required vs H100
0.989 PFLOPS peak). }(Analytical bound --- not a benchmark: these are
first-order model checks, not a measurement of actual serving
throughput.)*

\textbf{Worked example arithmetic details:}

\begin{itemize}
\tightlist
\item
  \textbf{Decode bandwidth.} In auto-regressive decode, generating one
  token requires a full forward pass of the 70B parameter matrix. At
  FP16 (2 bytes per parameter), the weight footprint is
  \(W = N \times 2 = 140\) GB. \textbf{This is the aggregate model-wide
  weight traffic} --- but note the unit of comparison carefully. For a
  hypothetical single-device implementation holding the entire FP16
  model, the aggregate model-weight traffic per decode step is
  \textasciitilde140 GB; on a tensor-parallel deployment both weight
  residency \emph{and} this traffic are sharded across GPUs, so the
  relevant comparison is per-GPU shard traffic against per-GPU HBM
  bandwidth, plus communication overhead. (A 140-GB FP16 model cannot
  reside on a single 80-GB H100 anyway, so the whole-model-on-one-GPU
  comparison is deliberately hypothetical.) With a TPOT budget of
  \textasciitilde25 ms per token, the \emph{aggregate} required
  bandwidth is
\end{itemize}

\[
B_\text{req} = \frac{W}{\tau} = \frac{140 \text{ GB}}{0.025 \text{ s}} = 5{,}600 \text{ GB/s} = 5.6 \text{ TB/s}
\]

{[}ILLUSTRATIVE{]}{[}DERIVED{]}. A single NVIDIA H100 HBM3 delivers
\textasciitilde3.35 TB/s peak bandwidth {[}1P{]}{[}FACT{]}. Since
\(5.6 > 3.35\), under this batch-1-equivalent weight-streaming model one
H100 cannot supply the required \emph{aggregate} weight-read rate ---
the decode phase is HBM-bandwidth-bound {[}ILLUSTRATIVE{]}{[}DERIVED{]}.
Read this as an \emph{aggregate lower-bound illustration}, not as the
literal bandwidth requirement of a single H100 in a realistic 8×H100
configuration (where the model and its traffic are sharded, so each GPU
streams only its shard against its own 3.35 TB/s, plus communication
overhead). This is a statement about the \emph{specified}
single-request, \textasciitilde25 ms/token model, not about all possible
70B serving: batching, weight-read amortization, and precision changes
(later chapters) fundamentally alter the per-request bandwidth demand.
Within the model, meeting that latency target calls for multi-GPU
scaling or bandwidth-increasing topologies (e.g.~NVLink-connected
nodes).

\begin{itemize}
\tightlist
\item
  \textbf{Prefill FLOPs.} The prefill pass computes attention over the
  9.2K input tokens and produces the first output token. The
  \emph{parameterized-linear} FLOP floor for a dense transformer forward
  pass is well approximated as
  \(2 \times \text{params} \times \text{tokens}\) (the factor of 2
  accounts for multiply-add per parameter per token). Call this what it
  is: \textbf{the parameterized-linear FLOP floor, \(2NL\), excluding
  explicit attention-score/value aggregation and other non-matmul
  overhead.} The full-attention component is \emph{not} captured by
  \(2NL\) --- it grows \textasciitilde quadratically with sequence
  length (Ch 22), so for a 9.2K-token prefill \(2NL\) is a lower bound
  on total compute, not the whole story. Thus:
\end{itemize}

\[
\text{prefill FLOPs} = 2 \times N \times L = 2 \times 70 \times 10^9 \times 9.2 \times 10^3 \approx 1.288 \times 10^{15} \approx 1.29 \text{ PFLOP}
\]

{[}ILLUSTRATIVE{]}{[}DERIVED{]}. An NVIDIA H100 delivers
\textasciitilde989 TFLOPS (BF16 dense tensor-core peak, without
sparsity) {[}1P{]}{[}FACT{]}, which is 0.989 PFLOPS. (Here and
throughout we quote the \emph{dense}, no-sparsity tensor-core peak for
any precision; NVIDIA's datasheet prints the 2× \emph{with-sparsity}
figure --- 1,979 TFLOPS for FP16/BF16 --- as its headline number, so a
dense-vs-sparse note is required even in vendor material.) The required
rate against the \textasciitilde1.08 s prefill budget is

\[
\text{rate} = \frac{1.29 \text{ PFLOP}}{1.08 \text{ s}} \approx 1.19 \text{ PFLOPS}
\]

Crossing \(1.19 > 0.989\) --- \textbf{even before accounting for
attention overhead, kernel inefficiency, and sub-100\% MFU, this
first-order lower bound already exceeds the H100's theoretical peak.}
Therefore one H100 cannot satisfy the stated prefill budget under these
assumptions; the prefill phase is compute-bound
{[}ILLUSTRATIVE{]}{[}DERIVED{]}. Note the direction of this argument:
exceeding the \emph{theoretical peak} proves the target is impossible on
a single H100 under the stated model, but it does \textbf{not} imply any
rate below the peak is achievable --- real kernels land well under MFU
100\%. Meeting the latency SLO therefore calls for multiple GPUs (model
parallelism or data parallelism) or more efficient attention
implementations.

\section{Measurement}\label{measurement}

How do we know whether a given workload is prefill-bound or decode-bound
in practice? Three practical measurements:

\begin{enumerate}
\def\labelenumi{\arabic{enumi}.}
\item
  \textbf{Divide the TTFT into retrieval + prefill.} Measure the
  retrieval latency (vector search, document fetch) separately from the
  prefill latency (prompt processing). In the canonical scenario,
  retrieval takes \textasciitilde120 ms, leaving \textasciitilde1.08 s
  for prefill on a 9.2K-token prompt. If prefill exceeds this budget,
  the bottleneck is compute or memory, not retrieval.
\item
  \textbf{Log TPOT against current sequence length, at controlled
  batch/concurrency.} Track the time-per-token as the generated sequence
  grows. This is a \emph{scaling curve}, not a binary verdict. An
  \textbf{increasing} TPOT as the sequence lengthens reveals growing
  KV-attention cost --- each new token attends over a longer KV
  sequence, and decode attention work and KV cache traffic generally
  grow with sequence length even though cached K/V projections are not
  recomputed. A relatively \textbf{flat} TPOT suggests weight streaming
  (a fixed per-step cost --- reading the full 140 GB weight matrix each
  step) remains dominant. TPOT does not \emph{have} to stay constant;
  batching, kernel choice, occupancy, and scheduling can make the
  observed curve complicated, so measure it rather than asserting a
  trend.
\item
  \textbf{Measure actual HBM utilization vs FLOP utilization.} Using GPU
  profilers (NVIDIA Nsight, vLLM stats), watch the fraction of peak HBM
  bandwidth consumed versus the fraction of peak FLOPS consumed. If HBM
  utilization is near 100\% while FLOPS utilization is low, the kernel
  is bandwidth-bound (typical of decode). If FLOPS utilization is near
  100\% while HBM utilization is low, the kernel is compute-bound
  (typical of prefill).
\end{enumerate}

These measurements cost nothing but a profiler attach and a few
representative requests --- and they prevent the architect from applying
the wrong optimization lever.

\section{Common Mistakes}\label{common-mistakes}

\begin{itemize}
\item
  \textbf{Treating prefill and decode as the same bottleneck.} A
  dangerous mistake is assuming that what slows prefill will also slow
  decode, or vice versa. Prefill FLOPs scale with input length; decode
  bandwidth is constant per token. Confusing the two leads to
  over-provisioning compute for a bandwidth-bound workload, or
  under-provisioning compute for a compute-bound one.
\item
  \textbf{Assuming one H100 suffices for 70B inference.} The arithmetic
  in Section 4 shows that a single H100 cannot meet the decode bandwidth
  demand (5.6 TB/s vs 3.35 TB/s) nor the prefill compute demand (1.29
  PFLOP vs 0.989 PFLOPS). Attributing 70B serveability to a single GPU
  ignores the opposite bottlenecks and produces an architecture that
  fails latency SLOs.
\item
  \textbf{Ignoring the input/output token asymmetry.} A workload with
  9.2K input tokens and 300 output tokens imposes a \emph{substantial
  prefill compute burden} despite producing relatively few output tokens
  --- do not infer serving cost or resource pressure from output-token
  count alone. The two legs of the request have different cost
  structures and different optimization paths: decode is repeated per
  output token and dominates \emph{wall-clock latency} (at the canonical
  \textasciitilde25 ms/token, 300 output tokens ≈ 7.5 s of decode vs a
  \textasciitilde1.08 s prefill budget), while prefill is a large
  one-shot burst of compute and dominates the \emph{compute} burden (the
  9.2K-input attention pass). Quoting only per-token latency or only
  throughput/s hides this distinction.
\item
  \textbf{Using manufacturer-peak bandwidth/FLOPS without mapping to
  real kernels.} H100 3.35 TB/s HBM3 is the theoretical peak; actual
  sustained bandwidth for weight-read kernels may be 60--80\% of peak.
  H100 989 TFLOPS BF16 dense is the tensor-core peak; actual matmul
  throughput depends on kernel fusion, batching, and precision. Citing
  raw numbers without kernel context is an anti-catalog error --- the
  number must be matched to the actual serving kernel in use.
\end{itemize}

\section{Architecture Consequence}\label{architecture-consequence}

The opposite bottlenecks of prefill and decode have immediate
architectural consequences:

\begin{itemize}
\item
  \textbf{Prefill is compute-bound → optimization levers:} kernel fusion
  (flash-attention, scaled-dot-product attention with fewer reads),
  model parallelism (splitting the 70B parameters across GPUs), more
  GPUs in parallel, or prefill-specific servers that dedicate hardware
  to the bursty preload phase. Quantization (FP8/BF16) does \emph{not}
  reduce the underlying algorithmic operation count (the usual 2 ×
  params × tokens FLOPs are approximately unchanged --- the matmuls
  still multiply the same matrices); rather it reduces the precision,
  lowers the bytes transferred/stored, and can raise the
  \emph{effective} accelerator throughput (8-bit tensor cores run at
  higher peak rate).
\item
  \textbf{Decode is bandwidth-bound → optimization levers:} continuous
  batching (vLLM, PagedAttention) to amortize weight reads across many
  tokens in flight, KV cache offloading to SSD or system memory, P/D
  (prefill/decode) disaggregation (separate pools of GPUs for prefill
  bursts vs.~decode steady state), and higher-bandwidth memory
  topologies (NVSwitch, HBM3e).
\end{itemize}

The key architectural decision this enables: \textbf{can we disaggregate
prefill and decode onto different hardware?} If prefill needs more FLOPS
and decode needs more bandwidth, a single homogeneous GPU pool is
suboptimal. A two-pool architecture --- a prefill cluster optimized for
compute (more GPUs, higher FLOPS) and a decode cluster optimized for
bandwidth (faster HBM, larger NVLink domains) --- is therefore an
architecture worth testing. Whether it actually needs fewer GPUs is not
established by the asymmetry alone: it depends on achievable
utilization, duplicated model residency, KV-transfer cost, interconnect
topology, workload shape, burstiness, and the hardware assigned to each
pool, all of which must be measured. (The asymmetry makes it a candidate
worth benchmarking, not a concluded win.) design. This is the central
theme of \hyperref[chap:11]{Chapter 11} (Serving) and Pattern 4 (Prefill/decode
disaggregation).

\section{What We Still Don't Know}\label{what-we-still-dont-know}

\begin{itemize}
\item
  \textbf{Sustained HBM bandwidth for weight-read kernels.} The 3.35
  TB/s H100 figure is a theoretical peak; real serving kernels (vLLM,
  TGI, transformer-engine) may achieve different sustained rates
  depending on kernel fusion, graph optimization, and batch size. --- a
  hypothesis to be measured on real hardware before it is cited.
\item
  \textbf{Effect of continuous batching on the 2 × params × tokens
  approximation.} Continuous batching amortizes weight reads and
  improves hardware utilization, scheduling efficiency, and aggregate
  goodput --- but it does \textbf{not} reuse one request's transformer
  activations to avoid the forward-pass work of another unrelated
  request, so it does not reduce the mathematical FLOP count per
  request. Compute that \emph{is} re-used across requests sharing a
  common prefix is a different mechanism: \textbf{prefix caching} (or KV
  reuse), which skips the re-prefill of a shared prefix and thereby does
  cut real forward-pass work for that portion. The magnitude of the
  prefix-cache hit rate at fleet scale is not yet pinned down. --- to be
  benchmarked with realistic concurrency patterns.
\item
  \textbf{Cross-technology bandwidth numbers.} HBM3e (next-generation)
  promises \textasciitilde6+ TB/s per GPU, and AMD MI300X promises
  \textasciitilde5.3 TB/s. How these compare to the 5.6 TB/s decode
  demand shifts the GPU count equation but does not change the
  fundamental bandwidth-bound vs compute-bound classification.
  {[}ILLUSTRATIVE{]}{[}DERIVED{]} --- vendor-published numbers, verify
  against the specific generation in use.
\item
  \textbf{Effect of quantization on decode bandwidth vs prefill FLOPs.}
  Int4 or FP8 quantization reduces the weight bytes (e.g., 70B at FP8 =
  70 GB instead of 140 GB), which directly lowers the HBM read demand in
  decode; it does \emph{not} reduce the algorithmic operation count of
  the matmuls (the 2 × params × tokens FLOPs are approximately
  unchanged), though lower precision can raise the \emph{effective}
  throughput of the tensor cores at which those ops run. The direction
  of the shift is clear (decode bandwidth improves; prefill ceiling
  rises in effective rate), but the precise trade-off point where decode
  becomes compute-bound rather than bandwidth-bound depends on the
  quantization scheme and hardware support. --- workload-dependent.
\end{itemize}

\section{End-of-Chapter Mini-Case}\label{end-of-chapter-mini-case}

The architect is still in the early design conversation about the
internal Q\&A tool described in \hyperref[chap:1]{Chapter 1}. The team has settled on a
70B-class dense model for accuracy, and they have a rough traffic
estimate: \textasciitilde2,000 registered employees, each roughly as
active as the canonical scenario (\textasciitilde10 requests/s average,
peaks \textasciitilde40 rps), with prompts of \textasciitilde1,200
tokens of internal documents plus \textasciitilde8K retrieved context
and \textasciitilde300 tokens of answer.

From the token layer alone, the architect can already state the hard
numbers: each request carries \textasciitilde9.2K input tokens and
\textasciitilde300 output tokens. At 10 rps average, the system must
sustain \textasciitilde92,000 input tokens/s in prefill and
\textasciitilde3,000 output tokens/s in decode. At 40 rps peak, those
numbers jump to \textasciitilde368,000 and \textasciitilde12,000
respectively. The architect can also state the two opposite bottlenecks:
prefill will be compute-bound (\textasciitilde1.29 PFLOP per request at
9.2K input), and decode will be bandwidth-bound (a \textasciitilde140
GB/token weight-read requiring \textasciitilde5.6 TB/s sustained over
each \textasciitilde25 ms decode step). The team's immediate
architectural choice is whether to build a single homogeneous GPU
cluster or to separate prefill and decode pools --- the differing
bottleneck profiles make disaggregation worth evaluating, though whether
it uses fewer GPUs must be measured rather than assumed (Ch 11).

Before passing this workload to \hyperref[chap:4]{Chapter 4} (workload anatomy), the
architect locks in one more number: the tokens-per-request split. This
is the currency every downstream chapter will price against, and it has
already been established as \textasciitilde9.2K input +
\textasciitilde300 output per request. The rest of the design --- model
fitting, memory budget, latency budget, cost --- can now proceed in
token-denominated terms.
